\section{Preliminaries and Problem Formulation}
\label{sec:prelim}

We study \emph{training-free} reinforcement learning (RL): a regime in which a single frozen multimodal speech/omni model is steered at inference time by a reward signal, with no gradient ever propagated into its parameters and no architectural surgery. In this paper the regime is realized concretely as \emph{verifiable-reward selection}---best-of-$N$, reranking, and candidate-card retrieval---over a fixed candidate pool, using a frozen bi-encoder to score candidates. This section fixes notation, states the invariant that distinguishes training-free RL from conventional alignment, specifies the single inference-time operator we actually use, introduces the reward-spread quantity that governs when selection helps, and specifies the reward class we admit. The objects defined here are reused verbatim in the Theory section.

\subsection{Training-free RL as constrained inference-time reward maximization}
\label{subsec:tfrl}

Let $g_\theta$ denote a pretrained speech/omni model with \emph{frozen} parameters $\theta$, and let $x$ be an input (an utterance, a prompt, or a paired context). The inference-time degrees of freedom---the only quantities we are permitted to vary---are collected into an \emph{action} $z \in \cZ$, where $\cZ$ is the selection action space made precise in \Cref{subsec:operator}. Executing an action produces an output that is scored by a bounded scalar \emph{reward} $\Rwd\colon \cZ \to [0,1]$ (we write $\Rwd(z)$ for the reward of the output realized by $z$ when $x$ is fixed). A pretrained \emph{base} (or prior) distribution $\qo \in \Delta(\cZ)$ over actions is induced by the frozen model and the default inference procedure. Training-free RL searches over distributions $\qz \in \Delta(\cZ)$ for one that raises expected reward while remaining close to the prior.

\begin{definition}[Training-free RL objective]
\label{def:objective}
Fix $x$, a base distribution $\qo$, a reward $\Rwd$, and a temperature $\beta > 0$. The KL-regularized inference-time objective is
\begin{equation}
\label{eq:Fobj}
\Fobj{\qz} \;=\; \Exp_{z \sim \qz}\!\big[\Rwd(z)\big] \;-\; \beta\,\KL{\qz}{\qo},
\end{equation}
and the training-free RL problem is $\qs \;=\; \arg\max_{\qz \in \Delta(\cZ)} \Fobj{\qz}$.
\end{definition}

The objective \eqref{eq:Fobj} is strictly concave in $\qz$ over the simplex, and its maximizer is the classical Gibbs/exponential tilt of the prior \citep{korbak2022kl}.

\begin{proposition}[Gibbs optimum]
\label{prop:gibbs}
The unique maximizer of \eqref{eq:Fobj} is
\begin{equation}
\label{eq:gibbs}
\qs(z) \;=\; \frac{1}{\Zpart}\,\qo(z)\,\exp\!\Big(\tfrac{1}{\beta}\,\Rwd(z)\Big),
\qquad
\Zpart \;=\; \Exp_{z \sim \qo}\!\Big[\exp\!\big(\tfrac{1}{\beta}\Rwd(z)\big)\Big],
\end{equation}
with optimal value $\Fobj{\qs} = \beta \log \Zpart$.
\end{proposition}

\begin{proof}[Proof sketch]
Adding a Lagrange multiplier for the normalization constraint and differentiating the concave functional \eqref{eq:Fobj} with respect to $\qz(z)$ yields $\Rwd(z) - \beta(\log \qz(z) - \log \qo(z)) - \beta = \text{const}$, i.e. $\log \qz(z) = \log \qo(z) + \Rwd(z)/\beta + \text{const}$; exponentiating and normalizing gives \eqref{eq:gibbs}. Substituting back collapses the objective to $\beta\log\Zpart$.
\end{proof}

Equation~\eqref{eq:gibbs} is the formal anchor of the paper: every concrete inference-time procedure we analyze---best-of-$N$, soft best-of-$N$, reranking, and candidate-card selection---is read as an \emph{estimator of, or a sampler from, the tilted target} $\qs$, and is therefore comparable through the lens of how well it approximates \eqref{eq:gibbs} at a given query budget. The temperature $\beta$ interpolates between the prior ($\beta\to\infty$, $\qs\to\qo$) and greedy reward maximization ($\beta\to 0^+$, $\qs\to$ the $\arg\max_z \Rwd$ point mass), recovering best-of-$N$ selection in the low-temperature limit \citep{beirami2024bon,verdun2025softbon,khalaf2025hedgetune}.

\begin{definition}[No-weight-update invariant]
\label{def:invariant}
A procedure is \emph{training-free} if, for every input and every query budget, the model parameters $\theta$ and the model's computational graph are held fixed; all adaptation is realized through the choice of $\qz$ over the inference-time action space $\cZ$. Formally, the map $\theta \mapsto \theta$ is the identity throughout, and only $\qz$ is optimized.
\end{definition}

\Cref{def:invariant} is the operational contract of this work. It excludes parameter-efficient fine-tuning and any backpropagation through $g_\theta$; it admits only \emph{selection, reweighting, and reranking} over a fixed candidate pool. Because $\qs$ depends on $\theta$ only through $\qo$ and on the task only through $\Rwd$, the frozen model's pretrained competence is \emph{activated}, never overwritten---the central premise the paper develops empirically and theoretically.

\subsection{The selection operator: a frozen bi-encoder ranking candidate cards}
\label{subsec:operator}

Every reproducible result in this paper is produced by a single inference-time operator: a frozen bi-encoder that ranks a finite pool of candidate \emph{cards} by cosine similarity under a verifiable reward. We make this precise; no other operator is instantiated in our experiments.

For a query $x$ we are given a finite \emph{candidate pool} $C(x)=\{y_1,\dots,y_N\}$ (e.g., tool/intent schema cards, or transcript hypotheses). The action space is the probability simplex over this pool, $\cZ = \Delta\big(C(x)\big)$; an action $z$ is a distribution over candidates, and a \emph{selection} action collapses to the vertex that returns a single chosen card.

\begin{definition}[Frozen bi-encoder selection operator]
\label{def:operator}
Let $E$ be a frozen bi-encoder (in our experiments \texttt{omni-embed-nemotron-3b}, a SentenceTransformer) that maps the query and each candidate card to $\ell_2$-normalized vectors $E(x)$ and $E(y_i)$. The operator scores each candidate by cosine similarity $s_i = \langle E(x),\,E(y_i)\rangle$ and induces the base distribution $\qo(y_i) \propto \exp(s_i/\beta)$ over the pool. Training-free RL reweights $\qo$ toward the verifiable reward via the Gibbs tilt \eqref{eq:gibbs}, and the realized \emph{selection} returns $\arg\max_i \Rwd(y_i)$ over the ranked pool (best-of-$N$), or the top-ranked card when the reward is used to construct the pool rather than to pick from it.
\end{definition}

Two properties of \Cref{def:operator} matter for the analysis. First, the operator is a \emph{single-sample deterministic readout} over a geometric object (points on the unit sphere): it adds no output entropy of its own and can only reweight candidates the pool already contains. Second, its usefulness is entirely governed by how much the reward \emph{varies across the pool}: if the reward is nearly constant over $C(x)$, selection cannot improve on the prior, whereas a reward that separates good from bad candidates lets selection concentrate mass on the reward maximizer. We make this dependence explicit next.

\begin{remark}[Generative operators are out of scope]
\label{rem:generative}
A co-equal generative operator---an autoregressive thinker--talker model that \emph{produces} rather than \emph{selects} candidates, with best-of-$N$ or reward-guided decoding tilting its sampler toward \eqref{eq:gibbs}---is a natural extension \citep{qwen3omni,mudgal2024controlled}. We do not instantiate one locally, and we make no claim about it; it is discussed only as future work. All results below concern the frozen bi-encoder selection operator of \Cref{def:operator}.
\end{remark}

\subsection{Reward spread}
\label{subsec:spread}

Selection helps exactly to the extent that the reward is \emph{not flat} over the candidate pool. We quantify this with the reward spread, the object that the qualitative lens in the Theory section (OSA-1) ties to the optimization gain.

\begin{definition}[Reward spread and optimization gain]
\label{def:spread}
Fix $x$, $\qo$, $\Rwd$, and $\beta$. The \emph{reward spread} $\spread$ is the range of $\Rwd$ over the support of the prior, $\spread \;=\; \max_{z\in\operatorname{supp}(\qo)}\Rwd(z)-\min_{z\in\operatorname{supp}(\qo)}\Rwd(z)$ (to leading order the $\qo$-weighted variance $\operatorname{Var}_{\qo}(\Rwd)$ plays the same role), and the \emph{optimization gain} of training-free RL over the prior is
\begin{equation}
\label{eq:gain}
\gain \;=\; \Fobj{\qs} - \Fobj{\qo} \;=\; \beta\log\Zpart \;-\; \Exp_{z\sim\qo}\!\big[\Rwd(z)\big] \;\ge\; 0,
\end{equation}
with equality if and only if $\Rwd$ is $\qo$-almost-surely constant.
\end{definition}

Equation~\eqref{eq:gain} is the free-energy expression for the value of the tilt: $\gain$ is nonnegative, vanishes exactly when the reward carries no spread, and grows as the reward separates candidates. The Theory section develops this as a qualitative lens (the Donsker--Varadhan / free-energy identity, used to explain---not to newly prove---why gains appear) and pins its two endpoints with strict, machine-checked lemmas: a flat reward yields zero gain (\texttt{flat\_no\_gain}), and a non-constant reward yields strictly positive gain (\texttt{gain\_pos\_of\_nonconstant}). Empirically this is the reading of our contrast: content/intent tasks expose high reward spread across candidate cards and therefore admit large selection gains, whereas the frozen model's paralinguistic read-out exposes little spread and selection yields no significant gain.

\subsection{Verifiable versus model-judged rewards}
\label{subsec:reward}

The validity of every guarantee in \Cref{subsec:tfrl,subsec:spread} rests on $\Rwd$ being trustworthy. We admit only \emph{verifiable} rewards.

\begin{definition}[Verifiable reward]
\label{def:verifiable}
A reward $\Rwd$ is \emph{verifiable} if it is a deterministic, label-derived function of the output and a ground-truth annotation $y^\star$, computable without invoking a learned scorer. Canonical instances are (i) negative word error rate $1-\mathrm{WER}(y,y^\star)$ for ASR/ST, (ii) exact-match/accuracy for classification and intent, and (iii) retrieval hit$@k$ / probe accuracy against a fixed label. A reward is \emph{model-judged} if it is the output of a separate generative or preference model evaluating $y$.
\end{definition}

\begin{assumption}[Verifiability]
\label{ass:verifiable}
Throughout, $\Rwd$ is verifiable in the sense of \Cref{def:verifiable}: bounded in $[0,1]$, deterministic given $y^\star$, and independent of $g_\theta$.
\end{assumption}

We adopt \Cref{ass:verifiable} for two reasons. First, model-judged rewards are susceptible to \emph{reward hacking}---the optimizer exploits gaps between the proxy and the true objective, so that increasing $\Exp_{\qs}[\Rwd]$ no longer increases task quality \citep{skalse2022rewardhacking,gao2023overopt}. Because the tilt \eqref{eq:gibbs} sharpens precisely along the direction of $\Rwd$, any reward misspecification is amplified rather than averaged out as $\beta\to 0$. Second, learned judges carry systematic artifacts---position bias and self-preference---that contaminate the candidate ranking \citep{positionbias2024,selfpreference2024}; verifiable rewards are invariant to these. Restricting to label-derived $\Rwd$ is therefore not a convenience but a precondition for the no-weight-update invariant of \Cref{def:invariant} to translate into genuine, contamination-free activation of pretrained competence \citep{tulu3}. With \Cref{def:objective,def:invariant,def:operator,def:spread,def:verifiable} and \Cref{ass:verifiable} in place, the notation required by the Theory section is complete.